\section*{第 9 章}
\addcontentsline{toc}{section}{第 9 章}

\exercise{习题 9.1}
\textbf{题目：}
求下列二元码字之间的汉明距离：
\[
\text{(1) }0000,\ 0101;\quad
\text{(2) }01110,\ 11100;
\]
\[
\text{(3) }010101,\ 101001;\quad
\text{(4) }1110111,\ 1101011.
\]

\par\textbf{解：}
逐位比较可得
\[
d_H(0000,0101)=2,
\]
\[
d_H(01110,11100)=2,
\]
\[
d_H(010101,101001)=4,
\]
\[
d_H(1110111,1101011)=3.
\]

\medskip
\exercise{习题 9.2}
\textbf{题目：}
某码字集合为
\[
\begin{array}{cccc}
0000000 & 1000111 & 0101011 & 0011101\\
1101100 & 1011010 & 0110110 & 1110001
\end{array}
\]
求最小汉明距离及纠错能力。

\par\textbf{解：}
两两比较可知任意两个不同码字之间的距离都为 $4$，因此
\[
d_{\min}=4.
\]
纠错能力为
\[
t=\left\lfloor \frac{d_{\min}-1}{2}\right\rfloor=\left\lfloor \frac{3}{2}\right\rfloor=1.
\]
故该码\textbf{可纠正 1 位错误}。

\medskip
\exercise{习题 9.3}
\textbf{题目：}
设二进制对称信道误码率 $p=10^{-2}$。
\begin{enumerate}
\item $(5,1)$ 重复码通过此信道传输时，不可纠正错误出现的概率是多少？
\item $(4,3)$ 偶校验码通过此信道传输时，不可检出错误出现的概率是多少？
\end{enumerate}

\par\textbf{解：}
对 $(5,1)$ 重复码，当 $5$ 位中出现 $3$ 位及以上错误时将不可纠正，因此
\[
P_e=\sum_{i=3}^{5}\binom{5}{i}p^i(1-p)^{5-i}
\approx 9.8506\times 10^{-6}.
\]

对 $(4,3)$ 偶校验码，出现偶数位错误而又非零时不能检出，因此
\[
P_u=\binom{4}{2}p^2(1-p)^2+p^4
\approx 5.88\times 10^{-4}.
\]

\medskip
\exercise{习题 9.4}
\textbf{题目：}
有一组等重码，每个码字长为 $5$，其中恰有 $3$ 个“1”。问该等重码是否为线性码？

\par\textbf{解：}
线性码必须包含全零码字。但该等重码中每个码字都恰有 $3$ 个“1”，不可能包含 $00000$。因此它\textbf{不是线性码}。

\medskip
\exercise{习题 9.5}
\textbf{题目：}
已知一个 $(7,4)$ 码的生成矩阵为
\[
G=
\begin{pmatrix}
1&0&0&0&1&1&1\\
0&1&0&0&1&0&1\\
0&0&1&0&0&1&1\\
0&0&0&1&1&1&0
\end{pmatrix},
\]
生成下列信息组的码字：
\[
(0100),\ (0101),\ (1110),\ (1001).
\]

\par\textbf{解：}
按 $c=mG$（模 2）计算，得
\[
(0100)\to 0100101,
\]
\[
(0101)\to 0101011,
\]
\[
(1110)\to 1110001,
\]
\[
(1001)\to 1001001.
\]

\medskip
\exercise{习题 9.6}
\textbf{题目：}
已知一个系统 $(7,4)$ 汉明码的监督矩阵为
\[
H=
\begin{pmatrix}
1&1&1&0&1&0&0\\
0&1&1&1&0&1&0\\
1&1&0&1&0&0&1
\end{pmatrix}.
\]
求：
\begin{enumerate}
\item 生成矩阵 $G$；
\item 当输入信息序列 $m=(110101101010)$ 时，求输出码序列。
\end{enumerate}

\par\textbf{解：}
将
\[
H=[P^T\ I_3],
\]
其中
\[
P^T=
\begin{pmatrix}
1&1&1&0\\
0&1&1&1\\
1&1&0&1
\end{pmatrix},
\]
故
\[
P=
\begin{pmatrix}
1&0&1\\
1&1&1\\
1&1&0\\
0&1&1
\end{pmatrix},
\qquad
G=[I_4\ P].
\]
即
\[
G=
\begin{pmatrix}
1&0&0&0&1&0&1\\
0&1&0&0&1&1&1\\
0&0&1&0&1&1&0\\
0&0&0&1&0&1&1
\end{pmatrix}.
\]

将输入序列按 4 位分组：
\[
1101,\quad 0110,\quad 1010.
\]
对应码字分别为
\[
1101001,\qquad 0110001,\qquad 1010011.
\]
故输出码序列为
\[
1101001\ 0110001\ 1010011.
\]

\medskip
\exercise{习题 9.7}
\textbf{题目：}
已知非系统码的生成矩阵为
\[
G'=
\begin{pmatrix}
0&0&0&1&0&1&1\\
0&0&1&0&1&1&0\\
0&1&0&1&1&0&0\\
1&0&1&1&0&0&0
\end{pmatrix}.
\]
求：
\begin{enumerate}
\item 等价系统码生成矩阵；
\item 典型监督矩阵。
\end{enumerate}

\par\textbf{解：}
对 $G'$ 作初等行变换，可化为系统形式
\[
G=
\begin{pmatrix}
1&0&0&0&1&0&1\\
0&1&0&0&1&1&1\\
0&0&1&0&1&1&0\\
0&0&0&1&0&1&1
\end{pmatrix}.
\]
因此
\[
P=
\begin{pmatrix}
1&0&1\\
1&1&1\\
1&1&0\\
0&1&1
\end{pmatrix},
\]
典型监督矩阵为
\[
H=
\begin{pmatrix}
1&1&1&0&1&0&0\\
0&1&1&1&0&1&0\\
1&1&0&1&0&0&1
\end{pmatrix}.
\]

\medskip
\exercise{习题 9.8}
\textbf{题目：}
已知某线性分组码生成矩阵为
\[
G'=
\begin{pmatrix}
0&0&1&1&1&0&1\\
0&1&0&0&1&1&1\\
1&0&0&1&1&1&0
\end{pmatrix}.
\]
求：
\begin{enumerate}
\item 系统码生成矩阵；
\item 典型监督矩阵 $H$；
\item 若译码器输入 $y=(0011111)$，求其伴随式；
\item 若译码器输入 $y=(1000101)$，求其伴随式。
\end{enumerate}

\par\textbf{解：}
对 $G'$ 作行变换，可得系统形式
\[
G=
\begin{pmatrix}
1&0&0&1&1&1&0\\
0&1&0&0&1&1&1\\
0&0&1&1&1&0&1
\end{pmatrix}
=[I_3\ P].
\]
于是
\[
P=
\begin{pmatrix}
1&1&1&0\\
0&1&1&1\\
1&1&0&1
\end{pmatrix}.
\]
故典型监督矩阵
\[
H=[P^T\ I_4]
=
\begin{pmatrix}
1&0&1&1&0&0&0\\
1&1&1&0&1&0&0\\
1&1&0&0&0&1&0\\
0&1&1&0&0&0&1
\end{pmatrix}.
\]

输入 $y=(0011111)$ 时，
\[
s=yH^T=(0010).
\]
输入 $y=(1000101)$ 时，
\[
s=yH^T=(1011).
\]

\medskip
\exercise{习题 9.9}
\textbf{题目：}
已知某 $(7,3)$ 码生成矩阵为
\[
G=
\begin{pmatrix}
1&0&0&1&0&1&1\\
0&1&0&1&1&1&0\\
0&0&1&0&1&1&1
\end{pmatrix}.
\]
求可纠正错误图样和对应伴随式。

\par\textbf{解：}
该矩阵已是系统形式 $G=[I_3\ P]$，因此
\[
H=[P^T\ I_4]
=
\begin{pmatrix}
1&1&0&1&0&0&0\\
0&1&1&0&1&0&0\\
1&1&1&0&0&1&0\\
1&0&1&0&0&0&1
\end{pmatrix}.
\]
伴随式共有 $2^4=16$ 种。若按\textbf{最小重量余集首部}选取可纠正错误图样，则可取：
\[
\begin{array}{c|c}
\text{伴随式} & \text{可纠正错误图样}\\
\hline
0000 & 0000000\\
0001 & 0000001\\
0010 & 0000010\\
0011 & 0000011\\
0100 & 0000100\\
0101 & 0000101\\
0110 & 0000110\\
0111 & 0010000\\
1000 & 0001000\\
1001 & 0001001\\
1010 & 0001010\\
1011 & 1000000\\
1100 & 0001100\\
1101 & 0001101\\
1110 & 0100000\\
1111 & 0011000
\end{array}
\]
据此即可构造标准阵列并完成译码。

\medskip
\exercise{习题 9.10}
\textbf{题目：}
题表中给出了 4 种 $(3,2)$ 码，判断各码组是否为线性码、是否为循环码。

\par\textbf{解：}
\[
\begin{array}{c|c|c}
\text{码组} & \text{是否线性} & \text{是否循环}\\
\hline
c_1 & \text{是} & \text{是}\\
c_2 & \text{是} & \text{否}\\
c_3 & \text{否} & \text{否}\\
c_4 & \text{否} & \text{否}
\end{array}
\]
其中 $c_1$ 满足加法封闭且满足循环移位封闭；$c_2$ 满足线性封闭，但不满足循环封闭；$c_3$ 不满足线性封闭；$c_4$ 连全零码字都不包含，因此既不是线性码，也不是循环码。

\medskip
\exercise{习题 9.11}
\textbf{题目：}
在 $\mathrm{GF}(2)$ 上计算：
\begin{enumerate}
\item $(x^4+x^3+x^2+1)+(x^3+x^2)$；
\item $(x^3+x^2+1)(x+1)$；
\item $x^4+x \pmod{x^2+1}$；
\item $x^4+1 \pmod{x+1}$；
\item $x^3+x^2+x+1 \pmod{x+1}$。
\end{enumerate}

\par\textbf{解：}
在 $\mathrm{GF}(2)$ 上同次项系数按模 2 相加，故
\[
\text{(1)}\quad (x^4+x^3+x^2+1)+(x^3+x^2)=x^4+1.
\]
\[
\text{(2)}\quad (x^3+x^2+1)(x+1)=x^4+x^2+x+1.
\]
\[
\text{(3)}\quad x^4+x \equiv x+1 \pmod{x^2+1}.
\]
\[
\text{(4)}\quad x^4+1 \equiv 0 \pmod{x+1}.
\]
\[
\text{(5)}\quad x^3+x^2+x+1 \equiv 0 \pmod{x+1}.
\]

\medskip
\exercise{习题 9.12}
\textbf{题目：}
在 $\mathrm{GF}(2)$ 域上计算
\[
x^7 \div g(x),\qquad g(x)=x^3+x+1,
\]
求其余式。

\par\textbf{解：}
多项式除法结果为
\[
x^7=(x^4+x^2+x+1)(x^3+x+1)+1.
\]
因此商式为
\[
q(x)=x^4+x^2+x+1,
\]
余式为
\[
r(x)=1.
\]
故
\[
x^7 \equiv 1 \pmod{x^3+x+1}.
\]

\medskip
\exercise{习题 9.13}
\textbf{题目：}
若已知一个 $(7,3)$ 循环码的生成多项式为
\[
g(x)=x^4+x^3+x^2+1,
\]
求生成矩阵。

\par\textbf{解：}
生成多项式对应码字为
\[
0011101.
\]
其循环移位仍为码字，可取 3 个线性无关码字：
\[
1110100,\qquad 0111010,\qquad 0011101.
\]
故可构成生成矩阵
\[
G=
\begin{pmatrix}
1&1&1&0&1&0&0\\
0&1&1&1&0&1&0\\
0&0&1&1&1&0&1
\end{pmatrix}.
\]

\medskip
\exercise{习题 9.14}
\textbf{题目：}
已知
\[
x^7+1=(x^3+x^2+1)(x^3+x+1)(x+1),
\]
写出：
\begin{enumerate}
\item $(7,3)$ 循环码的生成多项式；
\item $(7,4)$ 循环码的生成多项式；
\item $(7,6)$ 偶校验码的生成多项式；
\item $(7,1)$ 重复码的生成多项式。
\end{enumerate}

\par\textbf{解：}
循环码的生成多项式必须是 $x^7+1$ 的因式，且次数为 $n-k$。

对 $(7,3)$ 码，$\deg g=4$，可取
\[
g(x)=(x^3+x^2+1)(x+1)=x^4+x^2+x+1,
\]
或
\[
g(x)=(x^3+x+1)(x+1)=x^4+x^3+x^2+1.
\]

对 $(7,4)$ 码，$\deg g=3$，可取
\[
g(x)=x^3+x^2+1,
\qquad \text{或}\qquad
g(x)=x^3+x+1.
\]

对 $(7,6)$ 偶校验码，$\deg g=1$，故
\[
g(x)=x+1.
\]

对 $(7,1)$ 重复码，$\deg g=6$，故
\[
g(x)=\frac{x^7+1}{x+1}=x^6+x^5+x^4+x^3+x^2+x+1.
\]

\medskip
\exercise{习题 9.15}
\textbf{题目：}
\begin{enumerate}
\item 证明
\[
g(x)=x^{10}+x^8+x^5+x^4+x^2+x+1
\]
是 $(15,5)$ 循环码的生成多项式；
\item 若信息码多项式
\[
u(x)=x^4+x+1,
\]
写出对应系统码多项式。
\end{enumerate}

\par\textbf{解：}
对 $(15,5)$ 循环码，应有 $\deg g=15-5=10$，且 $g(x)$ 必须整除 $x^{15}+1$。直接验证可得
\[
x^{15}+1=(x^5+x^3+x+1)g(x),
\]
故题给 $g(x)$ 的确是该码的生成多项式。

系统编码时先作
\[
x^{10}u(x)=x^{14}+x^{11}+x^{10},
\]
再除以 $g(x)$。其余式为
\[
r(x)=x^8+x^7+x^6+x.
\]
故系统码多项式为
\[
c(x)=x^{10}u(x)+r(x)
=x^{14}+x^{11}+x^{10}+x^8+x^7+x^6+x.
\]
对应码组为
\[
100110111000010.
\]

\medskip
\exercise{习题 9.16}
\textbf{题目：}
已知 $(15,11)$ 循环码的生成多项式为
\[
g(x)=x^4+x+1,
\]
写出该码系统码生成矩阵 $G=(I\ Q)$。

\par\textbf{解：}
这是一个 $(15,11)$ 码，故生成矩阵有 $11$ 行。逐次计算
\[
x^{14},x^{13},\dots,x^4
\]
分别除以 $g(x)$ 后得到的余式，即可得到各行校验部分。最终系统生成矩阵可写为
\[
G=
\begin{pmatrix}
1&0&0&0&0&0&0&0&0&0&0&1&0&0&0\\
0&1&0&0&0&0&0&0&0&0&0&1&1&0&1\\
0&0&1&0&0&0&0&0&0&0&0&1&1&1&1\\
0&0&0&1&0&0&0&0&0&0&0&1&1&1&0\\
0&0&0&0&1&0&0&0&0&0&0&0&1&1&1\\
0&0&0&0&0&1&0&0&0&0&0&1&0&1&0\\
0&0&0&0&0&0&1&0&0&0&0&0&1&0&1\\
0&0&0&0&0&0&0&1&0&0&0&1&0&1&1\\
0&0&0&0&0&0&0&0&1&0&0&1&1&0&0\\
0&0&0&0&0&0&0&0&0&1&0&0&1&1&0\\
0&0&0&0&0&0&0&0&0&0&1&0&0&1&1
\end{pmatrix}.
\]

\medskip
\exercise{习题 9.17}
\textbf{题目：}
已知 $(15,7)$ 循环码的生成多项式为
\[
g(x)=x^8+x^7+x^6+x^4+1,
\]
若信息码组
\[
u=(0011001),
\]
写出系统码码组。

\par\textbf{解：}
系统编码时先作
\[
u(x)x^8,
\]
再除以 $g(x)$。计算得余式为
\[
r(x)\leftrightarrow 11110110.
\]
因此系统码码组为
\[
0011001\,11110110,
\]
即
\[
001100111110110.
\]

\medskip
\exercise{习题 9.18}
\textbf{题目：}
若 $(7,3)$ 循环码的生成多项式为
\[
g(x)=x^4+x^2+x+1.
\]
\begin{enumerate}
\item 若编码器输入是 $110$，写出系统编码码字；
\item 若译码器输入是 $1011011$，求其模 $g(x)$ 的伴随式，并给出译码结果。
\end{enumerate}

\par\textbf{解：}
该生成多项式对应的码字为
\[
0010111,
\]
循环移位后可得系统码中与输入 $110$ 对应的码字为
\[
1100101.
\]

对接收序列 $1011011$ 除以 $g(x)$，得伴随式
\[
s(x)=0111.
\]
该伴随式与错误图样
\[
0010000
\]
对应，因此纠错后码字为
\[
1001011.
\]

\medskip
\exercise{习题 9.19}
\textbf{题目：}
用八进制形式表示下列多项式：
\begin{enumerate}
\item $x^4+x^3+1$；
\item $x^5+x^2+1$；
\item $x^6+x^4+x^3+x+1$；
\item $x^9+x^5+x^3+x^2+1$。
\end{enumerate}

\par\textbf{解：}
先写成二进制系数串，再转为八进制：
\[
\text{(1)}\ x^4+x^3+1 \leftrightarrow 11001_2=31_8,
\]
\[
\text{(2)}\ x^5+x^2+1 \leftrightarrow 100101_2=45_8,
\]
\[
\text{(3)}\ x^6+x^4+x^3+x+1 \leftrightarrow 1011011_2=133_8,
\]
\[
\text{(4)}\ x^9+x^5+x^3+x^2+1 \leftrightarrow 1000101101_2=1055_8.
\]

\medskip
\exercise{习题 9.20}
\textbf{题目：}
将下列八进制数写成多项式：
\[
23,\quad 45,\quad 105,\quad 721.
\]

\par\textbf{解：}
\[
23_8=10011_2 \leftrightarrow x^4+x+1,
\]
\[
45_8=100101_2 \leftrightarrow x^5+x^2+1,
\]
\[
105_8=1000101_2 \leftrightarrow x^6+x^2+1,
\]
\[
721_8=111010001_2 \leftrightarrow x^8+x^7+x^6+x^4+1.
\]

\medskip
\exercise{习题 9.21}
\textbf{题目：}
某 $(3,1,3)$ 卷积码满足
\[
g_1=100,\qquad g_2=101,\qquad g_3=111,
\]
画出该码编码器。

\par\textbf{解：}
约束长度为 $K=3$，故编码器含 $2$ 个延时单元。按生成序列可得三路输出分别为
\[
c_1(k)=u(k-2),
\]
\[
c_2(k)=u(k)\oplus u(k-2),
\]
\[
c_3(k)=u(k)\oplus u(k-1)\oplus u(k-2).
\]
因此编码器可描述为：输入 $u(k)$ 依次经过两个延时单元，三路抽头分别按上式接到并串变换器输出。也即：
\begin{itemize}
\item 第一路取第二级寄存器输出；
\item 第二路取输入与第二级寄存器输出异或；
\item 第三路取输入、一级寄存器输出、二级寄存器输出三者异或。
\end{itemize}

\medskip
\exercise{习题 9.22}
\textbf{题目：}
已知一个 $(2,1,5)$ 卷积码
\[
g^{(1)}=(11101),\qquad g^{(2)}=(10011),
\]
求：
\begin{enumerate}
\item 编码器框图；
\item 生成多项式；
\item 若输入信息序列为 $11010001$，求输出码序列。
\end{enumerate}

\par\textbf{解：}
由码元连接可得两路生成多项式为
\[
g_1(x)=1+x+x^2+x^4,
\qquad
g_2(x)=1+x^3+x^4.
\]
故编码器由 $4$ 个延时单元组成，第一路取输入与延迟 $1,2,4$ 的模 2 和，第二路取输入与延迟 $3,4$ 的模 2 和。

当输入序列为 $11010001$ 时，两路输出分别为
\[
c_1=100000001101,
\qquad
c_2=110001100011.
\]
按并串变换合成为一路，得到
\[
c=11\ 01\ 00\ 00\ 00\ 01\ 01\ 00\ 10\ 10\ 01\ 11.
\]

\medskip
\exercise{习题 9.23}
\textbf{题目：}
已知一个 $(3,1,3)$ 卷积码
\[
g_1(x)=1+x+x^2,\qquad
g_2(x)=1+x+x^2,\qquad
g_3(x)=1+x^2.
\]
求：
\begin{enumerate}
\item 编码器框图；
\item 状态图、树图。
\end{enumerate}

\par\textbf{解：}
约束长度 $K=3$，因此有 $2$ 个延时单元。三路输出可写为
\[
c_1(k)=u(k)\oplus u(k-1)\oplus u(k-2),
\]
\[
c_2(k)=u(k)\oplus u(k-1)\oplus u(k-2),
\]
\[
c_3(k)=u(k)\oplus u(k-2).
\]
所以编码器是一个 $2$ 级移位寄存器，三路输出按上式抽头异或得到。

若取状态为当前寄存器内容 $(u_{k-1},u_{k-2})$，则状态转移可写为：
\[
00\xrightarrow{0/000}00,\qquad 00\xrightarrow{1/111}10,
\]
\[
10\xrightarrow{0/010}01,\qquad 10\xrightarrow{1/001}11,
\]
\[
01\xrightarrow{0/111}00,\qquad 01\xrightarrow{1/000}10,
\]
\[
11\xrightarrow{0/001}01,\qquad 11\xrightarrow{1/111}11.
\]
据此即可画出状态图与对应树图。

\medskip
\exercise{习题 9.24}
\textbf{题目：}
已知一个 $(2,1,3)$ 卷积码编码器结构如题图所示，写出 $g^{(1)},g^{(2)}$，并画出状态图、树图、网格图。

\par\textbf{解：}
由题图抽头可读出：
\[
g^{(1)}(x)=1+x^2,\qquad g^{(2)}(x)=x+x^2.
\]
取状态为两级寄存器内容 $(s_1,s_2)$，则
\[
c_1=u\oplus s_2,\qquad c_2=s_1\oplus s_2,
\qquad (s_1',s_2')=(u,s_1).
\]
于是状态转移关系为
\[
00\xrightarrow{0/00}00,\qquad 00\xrightarrow{1/10}10,
\]
\[
01\xrightarrow{0/11}00,\qquad 01\xrightarrow{1/01}10,
\]
\[
10\xrightarrow{0/01}01,\qquad 10\xrightarrow{1/11}11,
\]
\[
11\xrightarrow{0/10}01,\qquad 11\xrightarrow{1/00}11.
\]
据此即可直接作出该卷积码的状态图、树图和网格图。

\medskip
\exercise{习题 9.25}
\textbf{题目：}
已知一卷积编码器结构如题图所示，求：
\begin{enumerate}
\item $(n,k,K)$；
\item $g^{(1)},g^{(2)}$；
\item 若 $x=(10111)$，求输出 $c$。
\end{enumerate}

\par\textbf{解：}
由题图可知该编码器有两路输出、一路输入，且含 $3$ 级记忆，因此
\[
(n,k,K)=(2,1,4).
\]

按图中的抽头结构，可取
\[
g^{(1)}(x)=1,\qquad g^{(2)}(x)=1+x+x^3.
\]
即第一路为系统输出，第二路为输入与延迟 $1$、延迟 $3$ 的模 2 和。

设初始状态全零，则对输入
\[
x=(10111),
\]
两路输出分别为
\[
c_1=10111,\qquad c_2=11110.
\]
按并串方式合成为
\[
c=11\ 01\ 11\ 11\ 10.
\]

\medskip
\exercise{习题 9.26}
\textbf{题目：}
信源信息速率为 $9.6\ \text{kbit/s}$，通过一个 $1/2$ 码率卷积编码器后，用 4PSK 调制信号传输，载波频率为 $1\ \text{MHz}$，且 4PSK 采用滚降系数为 $1$ 的频谱成形。求：
\begin{enumerate}
\item 4PSK 的符号速率；
\item 信道中传输的 4PSK 信号功率谱示意。
\end{enumerate}

\par\textbf{解：}
卷积编码后比特率变为
\[
R_b=\frac{1}{1/2}\times 9.6=19.2\ \text{kbit/s}.
\]
4PSK 每个符号携带 $2$ bit，因此符号速率为
\[
R_s=\frac{R_b}{2}=9.6\ \text{kBaud}.
\]

滚降系数 $\alpha=1$ 时，升余弦频谱的带通信号总带宽为
\[
B=(1+\alpha)R_s=2R_s=19.2\ \text{kHz}.
\]
因此功率谱密度以
\[
f_c=1\ \text{MHz}
\]
为中心，对称分布，总展宽约为
\[
19.2\ \text{kHz}.
\]

\medskip
\exercise{习题 9.27}
\textbf{题目：}
某码率为 $1/2$、约束长度 $K=3$ 的卷积码网格图如题图所示。若接收序列为
\[
11\ 11\ 00\ 00\ 01\ 10\ 11\cdots,
\]
求：
\begin{enumerate}
\item 用维特比算法译码后的输出序列；
\item 状态转移图；
\item 生成多项式。
\end{enumerate}

\par\textbf{解：}
按题图进行维特比译码，到第 7 步得到的最大似然路径为
\[
11\ 01\ 00\ 10\ 01\ 10\ 11,
\]
对应的信息序列为
\[
1011100.
\]

由题图可读出状态转移规则：
\[
00\xrightarrow{0/00}00,\qquad 00\xrightarrow{1/11}10,
\]
\[
10\xrightarrow{0/00}01,\qquad 10\xrightarrow{1/10}11,
\]
\[
01\xrightarrow{0/11}00,\qquad 01\xrightarrow{1/01}10,
\]
\[
11\xrightarrow{0/10}01,\qquad 11\xrightarrow{1/01}11.
\]
据此即可画出状态转移图。

当输入单个 $1$ 时，两路输出分别为
\[
101,\qquad 111,
\]
故生成多项式为
\[
g_1(x)=1+x^2,\qquad g_2(x)=1+x+x^2.
\]

\medskip
